%=====================================================================
\chapter{String Matching}\label{ch:strings}
%=====================================================================
\locator{6}{Part VI --- Strings, Limits and Practice}

\begin{outcomes}
By the end of this chapter you will be able to:
\begin{tight}
  \item \textbf{Trace} the naive, Rabin--Karp and Knuth--Morris--Pratt
        algorithms on a given text and pattern \emph{(CLO-8)}.
  \item \textbf{Derive} the prefix function and \textbf{explain} why it gives
        the correct shift after a mismatch.
  \item \textbf{Prove} that KMP runs in $\bigTh{n+m}$ by an amortised argument.
  \item \textbf{Identify} the inputs on which the naive and Rabin--Karp
        algorithms degrade, and \textbf{explain} why KMP has none.
  \item \textbf{Choose} an algorithm for a stated matching workload.
\end{tight}
\end{outcomes}

\begin{prereq}
\chref{hashing} for Rabin--Karp --- it is a rolling hash and inherits every
caveat of \chref{hashing}. \chref{amortised} for the KMP running-time proof,
which is an amortised argument.
\end{prereq}

\begin{keyterms}
pattern $\bullet$ text $\bullet$ occurrence $\bullet$ naive matching $\bullet$
rolling hash $\bullet$ Rabin--Karp $\bullet$ spurious hit $\bullet$ prefix
function $\bullet$ border $\bullet$ Knuth--Morris--Pratt $\bullet$ shift
\end{keyterms}

\section{The Problem, and the Obvious Algorithm}

\begin{definitionbox}[String Matching]
Given a text $T$ of length $n$ and a pattern $P$ of length $m \le n$, find
every shift $s$ such that $T[s+1 \mathinner{\ldotp\ldotp} s+m] = P$.
\end{definitionbox}

\dsfig{\aoaalg{\algname{Naive-Match}$(T, P)$}{
\For{$s \gets 0$ \textbf{to} $n-m$}
  \State $i \gets 1$
  \While{$i \le m$ \textbf{and} $T[s+i] = P[i]$}
    \State $i \gets i + 1$
  \EndWhile
  \If{$i = m+1$}
    \State report an occurrence at $s$
      \Comment{then restart at $s+1$ anyway}
  \EndIf
\EndFor
}}{\algname{Naive-Match}. After a mismatch at shift $s$ it begins again at
$s+1$, re-reading characters it has already examined. That re-reading is the
whole of the $\bigO{nm}$.}{naivematch}

\begin{conceptbox}[What the naive algorithm forgets]
Suppose $P = \texttt{abcabcd}$ and the first mismatch at shift $s$ occurs at
position 7 --- so $T[s{+}1 \mathinner{\ldotp\ldotp} s{+}6] = \texttt{abcabc}$.

\smallskip
The naive algorithm now tries shift $s+1$, comparing $P[1] = \texttt{a}$ with
$T[s{+}2] = \texttt{b}$. But we \emph{already know} $T[s+2] = \texttt{b}$,
because we just read it. Shifts $s+1$ and $s+2$ are both hopeless and could
have been skipped without looking at the text at all.

\smallskip
\textbf{The information to skip them is in the pattern, not the text.} Both
better algorithms exploit this, in different ways: Rabin--Karp summarises the
window with a hash, and KMP precomputes exactly how far it may safely jump.
\end{conceptbox}

\section{Rabin--Karp}

\begin{conceptbox}[Compare numbers, not strings]
Treat each window of $m$ characters as a number in base $d$ modulo $q$. Then
comparing a window with the pattern is \emph{one} integer comparison rather
than up to $m$ character comparisons.

\smallskip
The hash of the next window is computed from the current one in $\bigO{1}$:
\[ h_{s+1} = \bigl(d\,(h_{s} - T[s{+}1]\,d^{\,m-1}) + T[s{+}m{+}1]\bigr)
   \bmod q, \]
--- drop the leading character, shift, add the new trailing one. That is why
it is a \term{rolling} hash.

\smallskip
\textbf{Unequal hashes prove a mismatch. Equal hashes prove nothing}, so every
hash match must be verified character by character. A verification that fails
is a \term{spurious hit}, and the running time is
\[ \bigO{n + m} \;+\; \bigO{m \times (\text{number of spurious hits})}. \]
\end{conceptbox}

\begin{alertbox}[Rabin--Karp inherits every caveat of \chref{hashing}]
Its expected $\bigO{n+m}$ assumes spurious hits are rare, which assumes the
hash scatters windows well. That is \chref{hashing}'s simple uniform hashing
assumption, with the same two ways to fail: a weak modulus by accident, or a
text chosen by an adversary who knows $q$.

\smallskip
The defence is also the same: a large modulus, and if the input is untrusted, a
modulus chosen at random at run time.
\end{alertbox}

\section{Knuth--Morris--Pratt}

\begin{definitionbox}[Prefix Function]
For a pattern $P$ of length $m$, the \term{prefix function}
$\pi[i]$ is the length of the longest proper prefix of
$P[1 \mathinner{\ldotp\ldotp} i]$ that is also a suffix of it.

\smallskip
For $P = \texttt{ababaca}$:
\begin{center}
\begin{tabular}{@{}lccccccc@{}}
$i$      & 1 & 2 & 3 & 4 & 5 & 6 & 7 \\
$P[i]$   & a & b & a & b & a & c & a \\
$\pi[i]$ & 0 & 0 & 1 & 2 & 3 & 0 & 1 \\
\end{tabular}
\end{center}
$\pi[5] = 3$ because \texttt{aba} is both a prefix and a suffix of
\texttt{ababa}.
\end{definitionbox}

\begin{conceptbox}[Why that is exactly the right shift]
Suppose the pattern has matched $k$ characters and then fails. The matched text
is $P[1 \mathinner{\ldotp\ldotp} k]$. Any shift that could still succeed must
align a prefix of $P$ with a \emph{suffix} of what has already matched ---
because that suffix is the only part of the text still under the pattern.

\smallskip
So the next viable shift aligns the longest proper prefix of
$P[1\mathinner{\ldotp\ldotp}k]$ that is also a suffix of it. That is
$\pi[k]$, by definition. Shifting to anything shorter would skip a possible
occurrence; shifting to anything longer is impossible.

\smallskip
\textbf{The text pointer never moves backwards.} That single property gives
the $\bigTh{n+m}$ bound, with no hypothesis about the input at all.
\end{conceptbox}

\dsfig{\aoaalg{\algname{Prefix-Function}$(P)$}{
\State $\pi[1] \gets 0$;\quad $k \gets 0$
\For{$i \gets 2$ \textbf{to} $m$}
  \While{$k > 0$ \textbf{and} $P[i] \ne P[k+1]$}
    \State $k \gets \pi[k]$
      \Comment{fall back to a shorter border}
  \EndWhile
  \If{$P[i] = P[k+1]$}
    \State $k \gets k+1$
  \EndIf
  \State $\pi[i] \gets k$
\EndFor
\State \Return $\pi$
}}{\algname{Prefix-Function}. It is KMP run against the pattern itself, and it
is $\bigTh{m}$ by the same amortised argument.}{prefixfn}

\dsfig{\aoaalg{\algname{KMP-Match}$(T, P)$}{
\State $\pi \gets \algname{Prefix-Function}(P)$;\quad $k \gets 0$
\For{$i \gets 1$ \textbf{to} $n$}
  \While{$k > 0$ \textbf{and} $T[i] \ne P[k+1]$}
    \State $k \gets \pi[k]$
      \Comment{$i$ never decreases}
  \EndWhile
  \If{$T[i] = P[k+1]$}
    \State $k \gets k+1$
  \EndIf
  \If{$k = m$}
    \State report an occurrence at $i - m$
    \State $k \gets \pi[k]$
  \EndIf
\EndFor
}}{\algname{KMP-Match}. The text index $i$ advances on every iteration and
never retreats, which is what makes the whole thing linear.}{kmp}

\begin{theorem}[KMP]\label{thm:kmp}
\algname{KMP-Match} runs in $\bigTh{n+m}$ time in the worst case.
\end{theorem}

\begin{proof}
An amortised argument, in the pattern of \chref{amortised}: take the potential
$\Phi = k$, the number of characters currently matched.

\smallskip
$\Phi$ begins at 0, never goes negative, and each iteration of the
\textbf{for} loop increases it by at most 1 (line 7). So over $n$ iterations
the total increase is at most $n$.

\smallskip
Each iteration of the inner \textbf{while} loop \emph{decreases} $\Phi$ by at
least 1, since $\pi[k] < k$ always. A quantity that rises by at most $n$ in
total and never goes below 0 can fall at most $n$ times. So the inner loop
executes at most $n$ times \emph{in total across the whole run} --- not $n$
times per iteration.

\smallskip
Hence the matching phase is $\bigO{n}$, and the same argument applied to
\algname{Prefix-Function} gives $\bigO{m}$. The $\bigOm{n+m}$ is immediate
since every character is read.
\end{proof}

\begin{conceptbox}[The same looseness as \algname{Build-Heap}]
The naive bound for KMP is ``$n$ iterations, each with a \textbf{while} loop
that could run $m$ times, so $\bigO{nm}$''. That is correct and badly loose,
for exactly the reason \chref{structures}'s \algname{Build-Heap} analysis was:
it charges every iteration the worst case, and the expensive iterations are
paid for by the cheap ones that preceded them.

\smallskip
Three appearances now --- \algname{Build-Heap}, the dynamic array, and KMP ---
of one idea. If a bound looks like (count) $\times$ (worst case) and the
expensive operations consume something the cheap ones built up, the bound is
probably loose.
\end{conceptbox}

\section{Measured}

\begin{worked}[First, Do They Agree?]
\begin{center}
\textbf{2{,}000 random texts and patterns over a 4-letter alphabet.
Disagreements: 0.}
\end{center}
All three return identical lists of occurrences. The naive version is the
oracle, as in \chref{bruteforce}.
\end{worked}

\begin{worked}[Ordinary Text, Where Nothing Happens]
Random text over a 4-letter alphabet, $m = 8$. Comparisons divided by $n$:

\rowcolors{2}{white}{lightbg}
\begin{center}
\begin{tabular}{@{}rrrrr@{}}
\toprule
\rowcolor{primary!15}
$n$ & naive & Rabin--Karp & KMP & naive (ms) \\
\midrule
1{,}000{,}000 & 1.333 & 0.000 & 1.250 &  6.7 \\
2{,}000{,}000 & 1.334 & 0.000 & 1.250 & 12.8 \\
8{,}000{,}000 & 1.333 & 0.000 & 1.250 & 53.7 \\
\bottomrule
\end{tabular}
\end{center}

\smallskip
\textbf{All three are linear, and the naive algorithm is barely worse than
KMP} --- 1.333 comparisons per character against 1.250. On random text a
mismatch arrives after about $1/(1-1/|\Sigma|)$ characters, so the naive
algorithm almost never re-reads anything.

\smallskip
\textbf{Rabin--Karp's column is 0.000} because the hash rejects every window in
one integer comparison and the character-comparison counter never moves.

\smallskip
\textbf{This is the table that makes people think the choice does not matter.}
It is also the table that any benchmark on English text would produce.
\end{worked}

\begin{worked}[The Text That Destroys the Naive Algorithm]
$T = \texttt{aaaa}\ldots\texttt{a}$ of length 200{,}000, and
$P = \texttt{aa}\ldots\texttt{ab}$ of length $m$. Every shift matches $m-1$
characters and then fails.

\rowcolors{2}{white}{lightbg}
\begin{center}
\begin{tabular}{@{}rrrrrr@{}}
\toprule
\rowcolor{primary!15}
$m$ & naive comparisons & $/(nm)$ & KMP comparisons & naive (ms) & KMP (ms) \\
\midrule
   10 &     1{,}999{,}910 & 1.0000 & 399{,}991 &   3.55 & 0.85 \\
  100 &    19{,}990{,}100 & 0.9995 & 399{,}901 &  28.98 & 1.10 \\
1{,}000 &   199{,}001{,}000 & 0.9950 & 399{,}001 & 279.22 & 0.80 \\
\bottomrule
\end{tabular}
\end{center}

\smallskip
\textbf{The naive count is exactly $nm$} --- the ratio is 1.0000, 0.9995,
0.9950. The worst case is not an upper bound being approached; it is being
attained.

\smallskip
\textbf{KMP's count is flat at about $2n$}, independent of $m$. The pattern got
a hundred times longer and KMP did the same work.

\smallskip
\textbf{At $m = 1{,}000$, KMP is 349$\times$ faster.} And the input is not
contrived: runs of repeated characters occur in DNA, in log files, in padded
records and in compressed streams.
\end{worked}

\begin{worked}[And the Input That Destroys Rabin--Karp]
Rabin--Karp's expected bound assumes spurious hits are rare. Shrink the modulus
and watch:

\rowcolors{2}{white}{lightbg}
\begin{center}
\begin{tabular}{@{}rrrr@{}}
\toprule
\rowcolor{primary!15}
modulus $q$ & spurious hits & comparisons & time (ms) \\
\midrule
             13 & 153{,}522 & 204{,}708 & 53.1 \\
          1{,}009 &   2{,}007 &   2{,}917 & 47.6 \\
      1{,}000{,}003 &       0 &       208 & 47.8 \\
1{,}000{,}000{,}007 &       0 &       208 & 47.8 \\
\bottomrule
\end{tabular}
\end{center}

\smallskip
\textbf{The verification count explodes a thousandfold and the clock barely
moves} --- 53.1 ms against 47.8. That is worth pausing on, because it is an
honest negative result: on \emph{random} text a spurious candidate is rejected
after about 1.3 characters, so a hundred thousand of them cost almost nothing.

\smallskip
\textbf{The $\bigO{nm}$ bound needs both many candidates \emph{and} expensive
verification.} Here is an input with both --- $T = \texttt{a}^{n}$,
$P = \texttt{a}^{m}$, so every window is a genuine match and every
verification runs the full $m$:

\rowcolors{2}{white}{lightbg}
\begin{center}
\begin{tabular}{@{}rrrrr@{}}
\toprule
\rowcolor{primary!15}
$m$ & RK comparisons & $/(nm)$ & RK (ms) & KMP (ms) \\
\midrule
   10 &     1{,}999{,}910 & 1.0000 &   8.48 & 2.02 \\
  100 &    19{,}990{,}100 & 0.9995 &  40.21 & 2.14 \\
1{,}000 &   199{,}001{,}000 & 0.9950 & 368.07 & 2.28 \\
\bottomrule
\end{tabular}
\end{center}

\smallskip
\textbf{Exactly $nm$ again, and KMP is 161$\times$ faster at $m = 1{,}000$.}

\smallskip
\textbf{The first attempt at this experiment used $P = \texttt{a}^{m-1}
\texttt{b}$} --- the naive algorithm's worst case --- and measured
\emph{zero} comparisons, because that pattern's hash differs from every
window's and Rabin--Karp rejected all of them instantly. The two algorithms
have different worst cases, and assuming they share one was the mistake.
\end{worked}

\begin{conceptbox}[Choosing]
\rowcolors{2}{white}{lightbg}
\begin{longtable}{@{}L{2.8cm}L{3.0cm}L{3.0cm}L{5.0cm}@{}}
\toprule
\rowcolor{primary!15}
\textbf{Algorithm} & \textbf{Typical} & \textbf{Worst case} &
\textbf{Use when} \\
\midrule
\endhead
naive & $\bigO{n}$ & $\bigO{nm}$ & $m$ is tiny, the alphabet is large, and the
  text is not repetitive \\
Rabin--Karp & $\bigO{n+m}$ & $\bigO{nm}$ & searching for \emph{many} patterns
  at once --- one pass, one hash set \\
KMP & $\bigTh{n+m}$ & $\bigTh{n+m}$ & a guarantee is needed, or the text is
  repetitive, or it is streaming \\
\bottomrule
\end{longtable}

\smallskip
\textbf{Rabin--Karp's real advantage is not speed} --- KMP beats it and has a
guarantee. It is that a \emph{set} of patterns can be matched in one pass by
putting all their hashes in a hash set, which neither of the others can do.
That generalises to the Rabin--Karp fingerprint used in rsync and in
plagiarism detection.

\smallskip
\textbf{KMP's real advantage} is that the text pointer never moves backwards,
so it works on a stream that cannot be rewound --- a socket, a pipe, a file too
large to hold.
\end{conceptbox}

\begin{pitfall}
\begin{tight}
  \item \textbf{Benchmarking only on English text.} Measured, the naive
        algorithm was within 7\% of KMP there and 349$\times$ worse on
        repetitive input.
  \item \textbf{Assuming two algorithms share a worst case.} The pattern that
        is worst for naive matching made Rabin--Karp do \emph{zero}
        comparisons.
  \item \textbf{Forgetting to verify a Rabin--Karp hash match.} Equal hashes do
        not mean equal strings. Skipping the check makes it fast and wrong.
  \item \textbf{A small or fixed modulus on untrusted input.} \chref{hashing}'s
        flooding attack, in a new costume.
  \item \textbf{Getting $\pi$ off by one.} It is the longest \emph{proper}
        prefix that is also a suffix; $\pi[i] < i$ always, and
        $\pi[i] = i$ would make the matching loop hang.
  \item \textbf{Believing KMP is $\bigO{nm}$} because of the nested
        \textbf{while}. The amortised argument is the whole point.
  \item \textbf{Reaching for KMP when \texttt{std::search} will do.} For one
        short pattern in a non-repetitive text, the library is fine and it is
        already written.
\end{tight}
\end{pitfall}

\begin{lab}[Three Algorithms and Two Traps]
Reproduces all five tables. Expect thirty minutes.

\begin{lstlisting}[language=C++]
// ch21_strings.cpp -- naive, Rabin-Karp and KMP, on text and on traps.
// Build: cl /O2 /EHsc /std:c++17 ch21_strings.cpp    (see Appendix A)

// --- 1. Naive matching: try every alignment ------------------------
// After a mismatch at offset s it restarts at s+1, re-reading
// characters it has already looked at. That re-reading is the O(nm).

// --- 2. Rabin-Karp: compare numbers, not strings -------------------
// A rolling hash of each window. Equal hashes mean a POSSIBLE match
// that must be verified; unequal hashes mean a certain mismatch,
// decided in one comparison. Expected O(n+m); the worst case is every
// window colliding, which is Chapter 13's problem in a new costume.
        if (hp == ht) {                       // candidate -- verify it
            size_t i = 0;
            while (i < m) { ++cmps; if (t[s+i] != p[i]) break; ++i; }
            if (i == m) hits.push_back(s);
            else if (spurious) ++*spurious;   // hash agreed, text did not
        }
        ht = (ht + mod - (high * (unsigned char)t[s]) % mod) % mod;
        ht = (ht * base + (unsigned char)t[s + m]) % mod;

// --- 3. KMP: never re-read a character of the text -----------------
// pi[i] is the length of the longest proper prefix of p[0..i] that is
// also a suffix of it. On a mismatch, that is exactly how far the
// pattern can be shifted without missing an occurrence -- so the text
// pointer only ever moves forward, giving O(n+m) with no hypothesis
// about the input at all.
static std::vector<size_t> prefix_function(const std::string& p) {
    std::vector<size_t> pi(p.size(), 0);
    for (size_t i = 1; i < p.size(); ++i) {
        size_t k = pi[i - 1];
        while (k > 0 && p[i] != p[k]) k = pi[k - 1];
        if (p[i] == p[k]) ++k;
        pi[i] = k;
    }
    return pi;
}

// --- 4. All three agree, on every input used below -----------------
// --- 5. Ordinary text, where nothing much happens ------------------
// This is the case that makes people think the choice does not matter.

// --- 6. The text that destroys the naive algorithm -----------------
// text = aaaa...a, pattern = aaa...ab. Every alignment matches m-1
// characters and then fails, so the naive version does n*m
// comparisons -- the worst case, reached exactly.

// --- 7. The input that destroys Rabin-Karp -------------------------
// Shrink the modulus and every window becomes a candidate that must
// be verified. This is Chapter 13's adversarial hashing again.

// --- 8. Rabin-Karp's actual worst case -----------------------------
// The pattern a^(m-1)b does NOT do it: its hash differs from every
// window of a^n, so Rabin-Karp rejects all of them in one comparison
// each and never verifies anything. (The first version of this
// program measured exactly zero comparisons, which was the clue.)
//
// The bound needs every window to be a CANDIDATE. Take p = a^m: now
// every window hashes equal, every verification runs the full m
// characters, and the total is n*m.
\end{lstlisting}

\textbf{What to notice.} Parts 7 and 8 together. Part 7 shows a thousandfold
rise in verifications that costs 11\% of the running time; part 8 shows the
input where the same bound actually bites. A measurement that moves a counter
dramatically and the clock barely at all is telling you the counter is not what
costs.

\smallskip
\textbf{Then try to break it.} Use KMP to search for \emph{all} occurrences of
10{,}000 different patterns in one text, and compare with a Rabin--Karp that
puts all 10{,}000 hashes in a \texttt{unordered\_set} and makes a single pass.
Predict first. KMP must run 10{,}000 times; Rabin--Karp runs once. This is the
workload where the algorithm with the worse guarantee is the right answer, and
it is why Rabin--Karp is still in use.
\end{lab}

\begin{checkpoint}
\begin{tightnum}
  \item Compute the prefix function for \texttt{aabaaab} and state the shift
        KMP uses after a mismatch at position 6 \emph{(CLO-8)}.
  \item Why does KMP's inner \textbf{while} loop not make it $\bigO{nm}$? Give
        the potential function.
  \item The pattern $\texttt{a}^{m-1}\texttt{b}$ is the naive algorithm's worst
        case and makes Rabin--Karp do no work at all. Explain both facts from
        the algorithms' structure.
\end{tightnum}
\end{checkpoint}

\begin{chaptersummary}
\begin{tight}
  \item Naive matching re-reads text after a mismatch, giving $\bigO{nm}$ worst
        case and near-linear behaviour on random text.
  \item Rabin--Karp compares rolling hashes: $\bigO{n+m}$ expected, with every
        hash match verified. It inherits all of \chref{hashing}'s assumptions.
  \item KMP's prefix function gives the exact safe shift, so the text pointer
        never moves backwards and the bound is $\bigTh{n+m}$ with no
        hypothesis about the input.
  \item KMP's bound is proved by an amortised argument with $\Phi = k$ --- the
        same looseness that \algname{Build-Heap} and the dynamic array had.
  \item Measured on random text: 1.333 comparisons per character for naive and
        1.250 for KMP. The choice looks irrelevant.
  \item Measured on $T = \texttt{a}^{n}$: the naive count was exactly $nm$ and
        KMP's was flat at $2n$. At $m = 1{,}000$ KMP was 349$\times$ faster.
  \item Measured: a weak Rabin--Karp modulus produced 153{,}522 spurious hits
        and cost only 11\% more time, because on random text each verification
        fails immediately. The $\bigO{nm}$ bound needs an input with both many
        candidates and expensive verification --- $P = \texttt{a}^{m}$ gives
        exactly $nm$, and KMP is 161$\times$ faster.
  \item Rabin--Karp's real advantage is matching many patterns in one pass;
        KMP's is working on a stream that cannot be rewound.
\end{tight}
\end{chaptersummary}

\begin{reviewq}
  \item \textbf{(MCQ)} KMP's worst-case running time is: (a)~$\bigTh{nm}$;
        (b)~$\bigTh{n\lg m}$; (c)~$\bigTh{n+m}$; (d)~$\bigTh{n^{2}}$.
  \item \textbf{(MCQ)} In Rabin--Karp, equal hash values mean:
        (a)~a confirmed occurrence; (b)~a candidate that must be verified;
        (c)~a certain mismatch; (d)~the modulus is too small.
  \item \textbf{(Short)} Define the prefix function and explain why it gives
        the correct shift after a mismatch \emph{(CLO-8)}.
  \item \textbf{(Short)} Give the amortised argument for KMP's linear bound.
  \item \textbf{(Short)} Name the input that defeats the naive algorithm and
        the input that defeats Rabin--Karp, and say why they are different.
  \item \textbf{(Applied)} Trace KMP matching \texttt{ababaca} in
        \texttt{abababacaba}, listing the value of $k$ after each text
        character.
  \item \textbf{(Applied)} You must scan a 50 GB log stream, which cannot be
        rewound, for any of 200 known error signatures. Choose an algorithm (or
        combination), justify it against this chapter's measurements, and say
        what you would check about the signatures first.
\end{reviewq}
